\begin{figure}[htbp]
\centering
\begin{tikzpicture}[font=\footnotesize]

% ---- price forecast variants ----
\node[fgout=3.5cm] (base) at (-4.9, 0)
  {\textbf{Base forecast}\\\(\hat{P}^{\mathrm{base}}\)\\with the \mbox{ENTSO-E} load forecast};
\node[fgout=3.5cm] (gen) at (0, 0)
  {\textbf{Forecast with OBTF inputs}\\\(\hat{P}^{\mathrm{gen}}\)\\with \(\hat{L}\), \(\hat{G}^{\mathrm{ren}}\), \(\widehat{RL}\) instead};
\node[fgout=3.5cm, fgvar] (wea) at (4.9, 0)
  {\textbf{Weather-free forecast}\\\(\hat{P}^{\mathrm{gen},-\mathrm{wea}}\)\\without raw weather};

% ---- evaluation row ----
\node[fgout=2.4cm, fgvar] (naive) at (-4.9,-2.6)
  {\textbf{Naive benchmarks}\\\(\mathrm{N}_{d-1}\), \(\mathrm{N}_{d-7}\)};
\node[fgeval=3.6cm] (eval) at (0,-2.6)
  {\textbf{Evaluation}\\RMSE, MAE, one-sided\\Giacomini--White test};
\node[fgdata=2.2cm] (real) at (4.9,-2.6)
  {\textbf{Realized}\\\textbf{day-ahead prices}};

\node[fgrq=6.4cm] (rq) at (0,-4.5)
  {\textbf{RQ2:} How do explicit load and renewable-generation forecasts
   affect the day-ahead price forecast?};

% ---- flows ----
\draw[fgflow] (base.south) -- ++(0,-0.55) -| ([xshift=-0.9cm]eval.north);
\draw[fgflow] (gen.south) -- (eval.north);
\draw[fgaux]  (wea.south) -- ++(0,-0.55) -| ([xshift=0.9cm]eval.north);
\draw[fgaux]  (naive.east) -- (eval.west);
\draw[fgflow] (real.west) -- (eval.east);
\draw[fgflow] (eval.south) -- (rq.north);

\end{tikzpicture}
\caption[Price forecast comparison for RQ2.]
{Design of RQ2. The base forecast uses price history, calendar
features, weather, and the \mbox{ENTSO-E} load forecast, while
\(\hat{P}^{\mathrm{gen}}\) replaces the latter with the OBTF load,
renewable-generation, and residual-load forecasts. A weather-free variant and two
naive forecasts serve as secondary comparisons. Notation as in
Figure~\ref{fig:system_overview}.}
\label{fig:rq3_price_experiments_design}
\end{figure}
